% TOP_PROBLEM: {"id": 3119, "title": "Obstacle number of planar graphs", "category_id": 3, "category": {"id": 3, "name": "graph_theory", "display_name": "Graph Theory", "description": "Problems involving graphs, networks, and their properties.", "slug": "graph-theory", "order_index": 3, "created_at": "2026-07-31T15:26:25.671Z"}, "status": "open", "problem_number": "OPG-37357", "set_id": 12, "statusQualification": "The source contains two questions; the example question and the universal boundedness question are retained separately without an open-status claim for each."}
% FIRST_POSTED: 2026-10-01
% ENTRY_KIND: statement_only
% UnsolvedMath ID 3119; source catalogue code OPG-37357
% !TeX program = lualatex
% Definition and sources only; this file is not an LLM solution attempt.
\documentclass[11pt,a4paper]{article}
\usepackage[margin=25mm]{geometry}
\usepackage{fontspec}
\setmainfont{lmroman10-regular.otf}[BoldFont=lmroman10-bold.otf,
 ItalicFont=lmroman10-italic.otf,BoldItalicFont=lmroman10-bolditalic.otf]
\usepackage{amsmath,amssymb,mathtools}
\usepackage{unicode-math}
\setmathfont{latinmodern-math.otf}
\AtBeginDocument{\let\setminus\smallsetminus}
\usepackage{xurl,hyperref}
\hypersetup{colorlinks=true,urlcolor=blue}
\setcounter{secnumdepth}{0}
\setlength{\parindent}{0pt}
\setlength{\parskip}{5pt}
\setlength{\emergencystretch}{3em}
\begin{document}
\section*{Obstacle number of planar graphs}\label{3119}
{\small UnsolvedMath ID 3119 · OPG-37357 · Graph Theory · OpenGarden}
\subsection{Definitions and mathematical statement}
For a finite simple graph $G$, an obstacle representation consists of distinct points $p_v\in\mathbb R^2$, one for each vertex, together with finitely many closed polygonal obstacles avoiding all these points. For distinct vertices $u,v$, require
\[
 uv\in E(G)\quad\Longleftrightarrow\quad
 (p_up_v)\cap\bigcup_{O\in\mathcal O}O=\varnothing,
\]
where $(p_up_v)$ is the open straight segment between their representing points. Thus obstacles block every nonedge and no edge. A polygonal obstacle is a filled simple polygon; arbitrarily many sides are permitted. Vertex points may be chosen in general position to eliminate collinearity degeneracies.

The obstacle number $\operatorname{obs}(G)$ is the minimum number of obstacles in such a representation, with zero allowed. A graph is planar if it has some drawing in the plane with no crossings between edge interiors. This planarity condition is a property of the abstract graph; an obstacle representation need not be a crossing-free planar drawing.

The source asks two questions: does a planar graph with obstacle number greater than one exist, and is there an absolute finite constant $C$ such that $\operatorname{obs}(G)\le C$ for every planar graph $G$? The second question concerns arbitrary graph order and arbitrary planar graphs, not a fixed embedding. The first asks for a single example and is logically distinct from proving that planar obstacle numbers are unbounded. The definition retains both original questions without treating their historical status as identical.
\subsection{Short English statement}
Determine whether planar graphs ever need more than one obstacle, and whether all planar graphs admit a uniformly bounded number of obstacles.
\subsection{Sources}
\begin{itemize}
\item[S1] ULAM.AI and the UnsolvedMath Contributors, supplied problems.json, record 3119 (OPG-37357), OpenGarden. Catalogue identity and source transcription; CC BY 4.0.
\url{https://huggingface.co/datasets/ulamai/UnsolvedMath}.
\item[S2] Open Problem Garden, Obstacle number of planar graphs. Original problem and discussion.
\url{http://www.openproblemgarden.org/op/obstacle_number_of_planar_graphs}.
\end{itemize}
\end{document}
